\subsection{The unstabilized semilocal one-zero algebra: explicit coordinates,
boundary maps, and information loss}
\label{subsec:unstabilized-semilocal-one-zero}

The audited source introduces $A_{\varnothing p}$ only as an
``intermediate semilocal crossed-product algebra.''  It does not specify its
space, action, completion, inclusion, quotient, or imprimitivity module.  The
two-place construction of Connes, Consani, and Marcolli does, however, supply
the exact seed: $\mathbb Q_p\times\mathbb R$, acted on diagonally by
$\{\pm p^n:n\in\mathbb Z\}$, with generic and one-zero strata
\cite[Section~8.2]{connesConsaniMarcolli2007}.  We now perform the finite-place
construction directly.  This gives a natural meaning to the missing symbol;
it is not an assertion that this meaning was already present in the local
source or is the unique possible enrichment of that symbol.

Let $P$ be a nonempty finite set of rational primes and put
\[
 S_P=P\sqcup\{\infty\},\qquad
 \mathbb A_{S_P}=\left(\prod_{q\in P}\mathbb Q_q\right)\times\mathbb R,
\]
\[
 \Gamma_P=\mathbb Z[1/\!\prod_{q\in P}q]^{\times}
 =\left\{\varepsilon\prod_{q\in P}q^{n_q}:
   \varepsilon\in\{\pm1\},\ n_q\in\mathbb Z\right\}.
\]
The group $\Gamma_P\cong(\mathbb Z/2)\times\mathbb Z^P$ acts by diagonal
multiplication.  For $x=((x_q)_{q\in P},y)\in\mathbb A_{S_P}$ define its
literal zero set
\[
 Z(x)=\{q\in P:x_q=0\}\cup
 \begin{cases}\{\infty\},&y=0,\\ \varnothing,&y\ne0.\end{cases}
\]
No support coordinate is identified with another one in this definition.
Set
\[
 W_P^{(1)}=\{x:\#Z(x)\leq1\},\qquad
 W_{\varnothing}=\{x:Z(x)=\varnothing\},\qquad
 W_v=\{x:Z(x)=\{v\}\}.
\]
The complement of $W_P^{(1)}$ is a finite union of closed coordinate
intersections, so $W_P^{(1)}$ is an invariant locally compact open subspace.

\begin{definition}[Natural unstabilized one-zero algebra]
\label{def:unstabilized-one-zero-algebra}
All crossed products in this subsection are full crossed products.  Define
\begin{align*}
 \mathfrak A_P^{(1)}&=C_0(W_P^{(1)})\rtimes\Gamma_P,\\
 \mathfrak I_P&=C_0(W_{\varnothing})\rtimes\Gamma_P,\\
 \mathfrak D_{v,P}&=C_0(W_v)\rtimes\Gamma_P,\qquad
 \mathfrak D_P=\bigoplus_{v\in S_P}\mathfrak D_{v,P}.
\end{align*}
For $p\in P$, the source symbol receives the explicit candidate
\[
 A_{\varnothing p}^{(P)}
 :=C_0(W_{\varnothing}\cup W_p)\rtimes\Gamma_P,
\]
with $A_{\varnothing}^{(P)}=\mathfrak I_P$ and
$A_p^{(P)}=\mathfrak D_{p,P}$.
\end{definition}

\begin{theorem}[The literal one-zero extension]
\label{thm:unstabilized-one-zero-extension}
Restriction to the closed one-zero boundary gives the short exact sequence
\begin{equation}
 0\longrightarrow\mathfrak I_P
 \longrightarrow\mathfrak A_P^{(1)}
 \longrightarrow\mathfrak D_P\longrightarrow0.
 \label{eq:unstabilized-one-zero-extension}
\end{equation}
For every $v\in S_P$, the invariant open subspace
$W_{\varnothing}\cup W_v$ gives a morphism from the $v$th edge extension
to \eqref{eq:unstabilized-one-zero-extension}, equal to the identity on
$\mathfrak I_P$ and to the summand inclusion on the quotient.
Full and reduced versions of all these algebras agree.
\end{theorem}

\begin{proof}
Inside $W_P^{(1)}$, the generic stratum is open and its complement is the
finite disjoint union $\bigsqcup_{v\in S_P}W_v$.  Restriction of functions
therefore gives
\[
 0\to C_0(W_{\varnothing})\to C_0(W_P^{(1)})
 \to\bigoplus_{v\in S_P}C_0(W_v)\to0.
\]
The maps are $\Gamma_P$-equivariant, and the full crossed-product functor is
exact.  This proves \eqref{eq:unstabilized-one-zero-extension}.  Removing all
boundary strata except $W_v$ leaves an invariant open subspace; extension by
zero gives the asserted middle and quotient maps and is literally the
identity on the common generic ideal.  Finally, $\Gamma_P$ is abelian-by-finite
and hence amenable.  Exel's approximation and regular-representation theorems
identify the full and reduced cross-sectional algebras in this case
\cite[Theorems~20.6--20.7]{exel2017}.
\end{proof}

We next retain every unit coordinate and calculate a finite-prime edge.  Put
\[
 K_P=\prod_{q\in P}\mathbb Z_q^{\times},\qquad
 \ell_q=\log q.
\]
For $p\in P$, let $K_{P\setminus p}=\prod_{q\in P\setminus\{p\}}
\mathbb Z_q^{\times}$, with the empty product understood as one point, and
define
\[
 \beta_p((b_q)_{q\ne p},s)
 =((p b_q)_{q\ne p},s+\ell_p).
\]
The finite-prime boundary carrier is the compact mapping torus
\begin{equation}
 B_{p,P}=(K_{P\setminus p}\times\mathbb R)/\langle\beta_p\rangle.
 \label{eq:multiplace-finite-mapping-torus}
\end{equation}

\begin{theorem}[Complete coordinates for the $P$-semilocal $p$-edge]
\label{thm:multiplace-finite-edge-coordinates}
The action of $\Gamma_P$ on $W_{\varnothing}\cup W_p$ is free and proper.
As a set, its orbit space is
\begin{equation}
 \mathcal S_{P,p}=(K_P\times\mathbb R)\sqcup B_{p,P}.
 \label{eq:multiplace-finite-edge-space}
\end{equation}
The symbol $\sqcup$ here records the two strata, not the topological-sum
topology; the cross-stratum convergence is specified below.
For a generic point, write
\[
 m_q=v_q(x_q),\qquad u_q=q^{-m_q}x_q\in\mathbb Z_q^{\times},
 \qquad\epsilon=\sgn(y).
\]
Its orbit coordinates are, without suppressing any cross-prime factor,
\begin{equation}
 a_q=\epsilon\!\prod_{r\in P\setminus\{q\}}r^{-m_r}u_q,
 \qquad
 t=\log|y|-\sum_{r\in P}m_r\ell_r.
 \label{eq:multiplace-generic-coordinates}
\end{equation}
For a point of $W_p$, define for $q\ne p$
\begin{equation}
 b_q=\epsilon\!\prod_{r\in P\setminus\{p,q\}}r^{-m_r}u_q,
 \qquad
 s=\log|y|-\sum_{r\ne p}m_r\ell_r;
 \label{eq:multiplace-p-boundary-coordinates}
\end{equation}
its orbit coordinate is $[(b_q)_{q\ne p},s]\in B_{p,P}$.

A generic net $(a^{(j)},t_j)$ converges to $[b,s]\in B_{p,P}$ exactly when
\begin{equation}
 t_j\longrightarrow-\infty,\qquad
 [(a_q^{(j)})_{q\ne p},t_j]\longrightarrow[b,s]
 \quad\hbox{in }B_{p,P};
 \label{eq:multiplace-finite-cross-convergence}
\end{equation}
there is no condition on the discarded coordinate
$a_p^{(j)}\in\mathbb Z_p^{\times}$.
\end{theorem}

\begin{proof}
Let $\gamma=\delta\prod_{r\in P}r^{n_r}\in\Gamma_P$.  On the generic
stratum,
\[
 m_q'=m_q+n_q,\qquad
 u_q'=\delta\!\prod_{r\ne q}r^{n_r}u_q,
 \qquad\epsilon'=\delta\epsilon.
\]
Substitution into \eqref{eq:multiplace-generic-coordinates} cancels every
$n_r$ and both factors $\delta$; the logarithmic coordinate is invariant
because $\log|\gamma|=\sum_rn_r\ell_r$.  Conversely, multiplication by the
unique element
\[
 \epsilon\prod_{r\in P}r^{-m_r}
\]
makes $y$ positive and every finite valuation zero.  Hence $(a,t)$ is a
complete generic orbit coordinate, not only an invariant.

On $W_p$, the same normalization is available for the sign and for all
valuations $m_q$ with $q\ne p$.  The unused power $p^k$ acts on the resulting
coordinates by $\beta_p^k$, which proves the boundary formula and its
uniqueness modulo \eqref{eq:multiplace-finite-mapping-torus}.

For the topology, let a generic representative have $p$-valuation $m_p$.
Equations \eqref{eq:multiplace-generic-coordinates} and
\eqref{eq:multiplace-p-boundary-coordinates} give the literal relations
\begin{equation}
 b_q=p^{m_p}a_q\quad(q\ne p),\qquad s=t+m_p\ell_p.
 \label{eq:multiplace-edge-coordinate-change}
\end{equation}
Thus $[(b_q),s]=[(a_q)_{q\ne p},t]$ in $B_{p,P}$.  Convergence of the
$p$-adic coordinate to zero is precisely $m_p\to+\infty$; if a boundary
representative remains in a compact set, the second identity forces
$t\to-\infty$ and gives \eqref{eq:multiplace-finite-cross-convergence}.
Conversely, use a local section of the proper $\mathbb Z$-quotient defining
$B_{p,P}$ to choose integers $m_j$ such that
\[
 \beta_p^{m_j}((a_q^{(j)})_{q\ne p},t_j)\longrightarrow(b,s).
\]
Because $t_j\to-\infty$ while the lifted real coordinates converge to $s$,
necessarily $m_j\to+\infty$.
Take all other valuations zero, set
$u_q^{(j)}=p^{m_j}a_q^{(j)}$ for $q\ne p$,
$x_p^{(j)}=p^{m_j}a_p^{(j)}$, and
$y_j=\exp(t_j+m_j\ell_p)$.  These are valid unit coordinates and converge to
the required boundary representative.  This proves both directions and the
absence of a condition on $a_p^{(j)}$.

The action is free because $y\ne0$.  For properness, let $C,D$ be compact
subsets of the edge and suppose $\gamma C\cap D\ne\varnothing$.  Every
$q\ne p$ coordinate is nonzero throughout the edge, so its valuations on
$C$ and $D$ bound $n_q$.  The real coordinate is also nonzero, so its absolute
values on $C,D$ bound $\sum_qn_q\ell_q$.  The already bounded $n_q$ for
$q\ne p$ then bound $n_p$, and the sign has only two values.  Only finitely
many $\gamma$ occur, proving properness.
\end{proof}

There is now an exact morphism to the one-ended local model, rather than a
Morita identification that erases the unit carrier.  Define
\[
 \rho_{p,P}:B_{p,P}\to S^1_{\ell_p},\qquad
 \rho_{p,P}([(b_q),s])=[s]_{\ell_p},
\]
and
\begin{equation}
 \Pi_{P,p}:\mathcal S_{P,p}\to X_{\ell_p},\qquad
 \Pi_{P,p}(a,t)=t,\quad
 \Pi_{P,p}([b,s])=[s]_{\ell_p}.
 \label{eq:multiplace-to-local-quotient}
\end{equation}
The mapping-torus relation changes $s$ by an integral multiple of $\ell_p$,
so $\rho_{p,P}$ is well defined.  The compact group $K_P$ acts by
\[
 k\cdot(a,t)=(ka,t),\qquad
 k\cdot[b,s]=[(k_qb_q)_{q\ne p},s].
\]
The second formula is well defined because coordinatewise multiplication by
$k$ commutes with $\beta_p$; the factor $k_p$ acts trivially on the boundary.
The convergence criterion proves that this action is continuous and that
$\Pi_{P,p}$ is exactly its orbit map.  Hence $\Pi_{P,p}$ is closed and proper.
Its generic fibres are $K_P$; its boundary fibres are copies of
$K_{P\setminus p}$.  Pullback therefore gives
\begin{equation}
\begin{array}{ccccccccc}
0&\to&C_0(\mathbb R)&\to&C_0(X_{\ell_p})&\to&C(S^1_{\ell_p})&\to&0\\
&&\downarrow\scriptstyle{\pi_P^*}&&\downarrow\scriptstyle{\Pi_{P,p}^*}
&&\downarrow\scriptstyle{\rho_{p,P}^*}\\
0&\to&C_0(K_P\times\mathbb R)&\to&C_0(\mathcal S_{P,p})
&\to&C(B_{p,P})&\to&0,
\end{array}
\label{eq:multiplace-finite-pullback-diagram}
\end{equation}
where $\pi_P(a,t)=t$.  Green imprimitivity, applied with the restricted
modules exactly as in the two-place proof, identifies the lower row with the
crossed-product edge extension
\[
 0\to\mathfrak I_P\to A_{\varnothing p}^{(P)}
 \to\mathfrak D_{p,P}\to0
\]
at the Morita level
\cite[Theorem~1.3.22]{meslandSengun2025}.

\begin{corollary}[Finite-place unit class with every unit fibre retained]
\label{cor:multiplace-finite-boundary-class}
Under
\[
 K_1(\mathfrak I_P)\cong
 K_1(C_0(K_P\times\mathbb R))
 \cong K_0(C(K_P))\cong C(K_P,\mathbb Z),
\]
the connecting map for the $p$-edge sends the Green image
$\eta_{p,P}$ of $[1_{B_{p,P}}]$ to
\begin{equation}
 \partial_{p,P}(\eta_{p,P})=-\mathbf1_{K_P}.
 \label{eq:multiplace-finite-boundary-class}
\end{equation}
\end{corollary}

\begin{proof}
The right vertical arrow of
\eqref{eq:multiplace-finite-pullback-diagram} is unital, so it sends the local
circle unit to the mapping-torus unit.  Theorem~\ref{thm:primitive-local-boundary}
gives $-1$ for the top boundary map.  Naturality sends it through
$\pi_P^*$, which is the constant-function inclusion
$\mathbb Z\to C(K_P,\mathbb Z)$.  This proves
\eqref{eq:multiplace-finite-boundary-class}.  The last $K_0$ identification
retains all locally constant integer-valued rank functions on the profinite
space $K_P$; it is not a normalization to one integer.
\end{proof}

The archimedean edge contains a different, equally explicit, information-loss
mechanism.  Normalize all finite valuations but do not choose a real sign.  If
$x_q=q^{m_q}u_q$, put
\begin{equation}
 c_q=\prod_{r\in P\setminus\{q\}}r^{-m_r}u_q,
 \qquad z=y\prod_{r\in P}r^{-m_r}.
 \label{eq:multiplace-arch-normalization}
\end{equation}
The residual sign acts by $(c,z)\mapsto(-c,-z)$.

\begin{theorem}[Archimedean edge and collapse of the two character labels]
\label{thm:multiplace-arch-boundary}
The action of $\Gamma_P$ on $W_{\varnothing}\cup W_\infty$ is free and
proper.  After Morita reduction of the free valuation directions, its
extension is
\begin{equation}
\begin{array}{ccccccccc}
0&\to&C_0(K_P\times\mathbb R^\times)\rtimes C_2
&\to&C_0(K_P\times\mathbb R)\rtimes C_2
&\to&C(K_P)\rtimes C_2&\to&0,
\end{array}
\label{eq:multiplace-arch-crossed-extension}
\end{equation}
where $C_2$ acts diagonally by $(c,z)\mapsto(-c,-z)$ and by
$c\mapsto-c$ on the boundary.  Constant functions in $K_P$ give a morphism
from the local reflection extension
\eqref{eq:primitive-archimedean-extension} to
\eqref{eq:multiplace-arch-crossed-extension}.  If
\[
 e_\pm=\frac{1\pm U}{2}\in C(K_P)\rtimes C_2,
\]
then
\begin{equation}
 \partial_{\infty,P}[e_+]
 =\partial_{\infty,P}[e_-]
 =-\mathbf1_{K_P}\in C(K_P,\mathbb Z).
 \label{eq:multiplace-arch-boundary-classes}
\end{equation}
Moreover $[e_+]=[e_-]$ in $K_0(C(K_P)\rtimes C_2)$.
\end{theorem}

\begin{proof}
Formula \eqref{eq:multiplace-arch-normalization} is invariant under every
prime-power direction, and every orbit has a unique representative with all
finite valuations zero, up to the residual sign.  This gives
\eqref{eq:multiplace-arch-crossed-extension}.  Freeness follows because all
finite coordinates are nonzero on this edge.  Properness follows because
their valuations on two compact sets bound every prime exponent; the residual
sign is finite.

The projection $K_P\times\mathbb R\to\mathbb R$ is equivariant for the
diagonal and reflection actions.  Pullback, followed by crossed product, is
therefore the promised extension morphism.  It sends the local projections
$p_\pm$ to $e_\pm$ and sends the local ideal generator to the constant rank
function.  Equation \eqref{eq:multiplace-arch-boundary-classes} follows from
Theorem~\ref{thm:primitive-archimedean-boundary} by naturality.

It remains to prove, rather than assume, the equality of the two boundary
classes.  Choose one $q_0\in P$.  If $q_0$ is odd, choose one representative
from every pair $\{r,-r\}$ in $(\mathbb Z/q_0\mathbb Z)^\times$.  If $q_0=2$,
choose the residue class $1\pmod 4$; its negative is the class
$3\pmod 4$.  The inverse image of the chosen classes defines a clopen set
$H\subset K_P$ with
\[
 K_P=H\sqcup(-H).
\]
Thus the principal double cover $K_P\to K_P/\{\pm1\}$ has the explicit
continuous section whose image is $H$.  Consequently
\[
 C(K_P)\rtimes C_2\cong M_2(C(H)),
\]
and $e_+$ and $e_-$ are rank-one projections in the same $2\times2$ matrix
algebra.  In the standard two-sheet matrix coordinates the constant unitary
$\operatorname{diag}(1,-1)$ conjugates $e_+$ to $e_-$, proving
$[e_+]=[e_-]$.  This is the exact point at which the natural semilocal carrier
forgets the local sign-character difference.
\end{proof}

We can now compute the natural multi-boundary map without pretending that the
whole action is proper.  Let
\[
 \Lambda_P=\bigoplus_{p\in P}\mathbb Z\epsilon_p
 \oplus\mathbb Z\epsilon_+\oplus\mathbb Z\epsilon_-.
\]
Define the boundary-label morphism
\begin{equation}
 R_P:\Lambda_P\longrightarrow K_0(\mathfrak D_P),\qquad
 R_P(\epsilon_p)=\eta_{p,P},\quad
 R_P(\epsilon_\pm)=[e_\pm],
 \label{eq:unstabilized-boundary-label-morphism}
\end{equation}
where the last two classes are transported through the valuation Morita
module.  Also define the exact constant-coordinate morphism
\[
 c_P:\mathbb Z\longrightarrow C(K_P,\mathbb Z),\qquad
 c_P(n)=n\mathbf1_{K_P}.
\]

\begin{theorem}[Unstabilized arithmetic realization of the codiagonal shadow]
\label{thm:unstabilized-codiagonal-shadow}
For the connecting map of
\eqref{eq:unstabilized-one-zero-extension},
\begin{equation}
 \partial_P^{(1)}R_P
 ( (k_p)_{p\in P},n_+,n_- )
 =-\left(\sum_{p\in P}k_p+n_++n_-\right)\mathbf1_{K_P}.
 \label{eq:unstabilized-codiagonal-formula}
\end{equation}
Equivalently, the square
\begin{equation}
\begin{CD}
 \Lambda_P @>{\delta_P^{\mathrm{bal}}}>> \mathbb Z\\
 @V{R_P}VV @VV{c_P}V\\
 K_0(\mathfrak D_P) @>{\partial_P^{(1)}}>> K_1(\mathfrak I_P)
\end{CD}
\label{eq:stable-to-unstabilized-k-square}
\end{equation}
commutes under the displayed Morita and suspension coordinates.

The morphism $R_P$ does not preserve all formal labels injectively:
\begin{equation}
 \ker R_P=\mathbb Z(\epsilon_+-\epsilon_-).
 \label{eq:unstabilized-boundary-label-kernel}
\end{equation}
Hence the formal scalar kernel $A_{|P|+1}$ maps onto the actual kernel on the
distinguished boundary-unit span, with exact kernel
$\mathbb Z(\epsilon_+-\epsilon_-)$.  The latter actual kernel has rank $|P|$
and an explicit $A_{|P|}$ basis after replacing the two archimedean labels by
their common image.
\end{theorem}

\begin{proof}
For each $v$, the edge algebra is an invariant open ideal in
$\mathfrak A_P^{(1)}$, and Theorem~\ref{thm:unstabilized-one-zero-extension}
gives a morphism of its extension into the global one-zero extension.  The
connecting maps are natural
\cite[Definition~21.1.1 and Example~21.1.2(a)]{blackadar1986}.  Therefore
Corollary~\ref{cor:multiplace-finite-boundary-class} supplies the term
$-k_p\mathbf1_{K_P}$ for every finite $p$, while
Theorem~\ref{thm:multiplace-arch-boundary} supplies
$-n_+\mathbf1_{K_P}$ and $-n_-\mathbf1_{K_P}$.  Additivity proves
\eqref{eq:unstabilized-codiagonal-formula} and the square.

The quotient $\mathfrak D_P$ is a direct sum over the distinct boundary
strata.  Each finite unit class has nonzero rank in its own summand and is
therefore independent of all other finite and archimedean summands.  In the
archimedean summand, the preceding theorem proves that the two classes agree;
their common rank is one, so this is their only relation.  This proves
\eqref{eq:unstabilized-boundary-label-kernel}.  The formal kernel of
$\delta_P^{\mathrm{bal}}$ is the already proved root lattice
$A_{|P|+1}$.  Quotienting by the displayed simple difference identifies the
two archimedean coordinates and leaves
\[
 \left\{((k_p)_{p\in P},n)\in\mathbb Z^{P}\oplus\mathbb Z:
       \sum_pk_p+n=0\right\},
\]
whose consecutive-difference basis is $A_{|P|}$.  No lattice coordinate has
been silently deleted: the lost coordinate and the map that loses it are
both displayed.
\end{proof}

There is a final analytic distinction.  Although every individual edge action
is free and proper, the simultaneous action on $W_P^{(1)}$ is not proper.

\begin{proposition}[Explicit obstruction to a global Green reduction]
\label{prop:unstabilized-global-action-nonproper}
For every nonempty $P$, the action of $\Gamma_P$ on $W_P^{(1)}$ is not proper.
Consequently one may not replace
\eqref{eq:unstabilized-one-zero-extension} by the commutative orbit-space
extension using Green imprimitivity.  This does not affect its crossed-product
exactness or the edgewise boundary calculation above.
\end{proposition}

\begin{proof}
Fix $p\in P$.  Choose integers $n_j\to\infty$ along a subsequence for which
the tuple $(p^{n_j})_{q\in P\setminus\{p\}}$ converges in the compact group
$K_{P\setminus p}$.  Put
\[
 z_j=((x_{q,j})_{q\in P},y_j),\qquad
 x_{p,j}=p^{n_j},\quad x_{q,j}=p^{n_j}\ (q\ne p),\quad y_j=1.
\]
Then $z_j$ converges to a point of the $p$-boundary.  For
$\gamma_j=p^{-n_j}$,
\[
 \gamma_jz_j=((1)_{q\in P},p^{-n_j}),
\]
which converges to a point of the archimedean boundary.  The union of both
sequences and their limits is compact in $W_P^{(1)}$, while the distinct
$\gamma_j$ escape every finite subset of the discrete group $\Gamma_P$.
Thus the inverse image of a compact subset of
$W_P^{(1)}\times W_P^{(1)}$ under $(\gamma,z)\mapsto(\gamma z,z)$ is not
compact.  This is exactly failure of properness.
\end{proof}

\paragraph{Exact comparison boundary before the extension lift.}
The natural algebra \eqref{eq:unstabilized-one-zero-extension} is now an
independently defined unstabilized arithmetic representative, and
\eqref{eq:stable-to-unstabilized-k-square} is an exact morphism from the
stable scalar boundary calculation to its unit-retaining $K$-theory shadow.
It is not an isomorphism: the explicit morphism
$R_P$ has kernel \eqref{eq:unstabilized-boundary-label-kernel}, while the ideal
map $c_P$ retains only the constant subgroup of the much larger group
$C(K_P,\mathbb Z)$.  This proves the precise information loss; it does not
assert that the two constructions are categorically unrelated.

Subsection \ref{subsec:stable-to-natural-extension} resolves that narrower
problem after stabilization.  It writes the ideal, quotient, middle,
multiplier, and corona maps explicitly, proves the Busby equation, and gives a
minimal symmetric split enrichment when the two sign classes must remain
distinct in $K_0$.  The nonconstant unit-fibre classes remain outside the
stable scalar image and are retained as such.  No statement about an
individual zeta zero follows from the construction.

\paragraph{Bounded coordinate replay.}
The script
\begin{center}
\footnotesize
\path{certificates/globalization_nte_multiplace_unstabilized/}\\
\path{verify_multiplace_unstabilized.py}
\end{center}
replays 362 named checks for one through eight finite primes.  They cover
exponents and signs, screw coordinates, the lattice square, the clopen
section, and the nonproperness sequence.  The accompanying JSON receipt has
stable ID
\begin{center}
\footnotesize
\texttt{CERT-GNTE-MULTIPLACE-}\\
\texttt{UNSTABILIZED-20260830-0001}.
\end{center}
It records two Python~3.13/SymPy runs with exit code zero.  It does not replace
the printed topology, crossed-product, Green-imprimitivity, or six-term
naturality proofs, and it launches no Lean process.
